\documentclass[12pt, titlepage]{article}
\usepackage[letterpaper, width = 150mm, top = 20mm, bottom = 25mm]{geometry}
\usepackage{authblk}
\usepackage{fancyhdr}
%\pagestyle{fancy}
\makeatother
\renewcommand{\headrulewidth}{0pt}
\renewcommand{\footrulewidth}{0pt}
\makeatletter
\usepackage[T1]{fontenc}
\usepackage{roboto}
\usepackage{nimbusserif}
\usepackage{graphicx}
\usepackage{amsmath, amsfonts}
\newcommand{\source}[1]{\footnotesize\caption*{\emph{Source}: {#1}}} %
\newcommand{\note}[1]{\footnotesize\caption*{\emph{Note}: {#1}}} %
\usepackage{booktabs}
\usepackage{graphicx}
\usepackage{adjustbox}
\usepackage[flushleft]{threeparttable}
\usepackage{authblk}
\usepackage[font=sf]{caption}
\usepackage{subcaption}
\usepackage{textpos}
\usepackage{rotating}
\usepackage{multirow}

\usepackage{ragged2e}
\usepackage{titlesec}
\usepackage[authoryear]{natbib}
\usepackage[colorlinks]{hyperref}
\hypersetup{
	linkcolor=blue,
	citecolor=blue,
	urlcolor=blue
}
\usepackage{relsize}
\usepackage{setspace}
\clubpenalty=100000
\widowpenalty=10000
\singlespacing
\pagenumbering{arabic}
\captionsetup{singlelinecheck=false}
\numberwithin{equation}{section}
\usepackage[hang,flushmargin, bottom]{footmisc}

\newcommand\blfootnote[1]{%
	\begingroup
	\renewcommand\thefootnote{}\footnote{#1}%
	\addtocounter{footnote}{-1}%
	\endgroup
}
% to change titles font family
\usepackage{titling}
%\usepackage[font=sf]{floatrow}
\usepackage{float}
\usepackage{enumitem}
\setlist{nolistsep}
%\floatsetup[table]{capposition=top}
\usepackage{sectsty}
\sectionfont{\sffamily\mdseries}
\subsectionfont{\sffamily\mdseries}
\subsubsectionfont{\sffamily\mdseries}
%\titleformat*{\section}{\Large\sffamily\MakeUppercase\clearpage}
%\titleformat*{\subsection}{\large\sffamily\MakeUppercase}
%\titleformat*{\subsubsection}{\normalsize\sffamily\MakeUppercase}


\renewcommand{\bibsection}{\section{References}}

\date{}

\renewcommand\Affilfont{\itshape\small}

\begin{document}
	\begin{titlepage}
		\centering
		%\includegraphics[width =4cm]{logic-2020.png}\\
		\LARGE{\sffamily{Growth, Debt, and Inequality}}\\
		\vspace{5mm}
		\normalsize{Zhuo Fu Chen\(^\ast\), Anthony Perez Eisenbarth\(^\dagger\),
			and Richard Atta Twumasi\(^\ddagger\)}\\ \vspace{5mm}
		\textit{
			\vspace{1mm}
			\(^\ast\)zhuofuchen@gmail.com, \(^\dagger\)ageisenbarth@gmail.com, \(^\ddagger\)richard.twumasi@outlook.com
		}
		\begin{Center}
			\textbf{Abstract}
		\end{Center}
		\justifying
		Standard macroeconomic theory assumes that persistent fiscal deficits
		lead to negative impacts on distributional outcomes. Though there is little empirical evidence to substantiate the assumption of this relationship, it nevertheless perseveres, guiding the study
		of both fiscal and monetary policy. The most recent global financial crisis
		has led to an unprecedented increase in public debt across the world, raising
		serious concerns about its economic impact. In our sample of 137 countries over a period 1960 to 2022, we study the relationship between debt, deficits, and inequality.
	\end{titlepage}
%	\tableofcontents
%	\addtocontents{toc}{\protect\thispagestyle{empty}}
%	\pagenumbering{gobble}
%	\cleardoublepage
	\pagenumbering{arabic}

\section{Introduction}

In the aftermath of the Great Recession and the COVID-19 pandemic, public debt levels have soared to unprecedented levels. The extraordinarily high debt levels and seeming fragility of the global economy has brought public and academic debate to reconsider the long-run implications of fiscal spending. Modern arguments can be traced back to \cite{Buchanan1958}, \cite{Meade1958}, and \cite{Modigliani1961}, although the issue goes back as far as Adam Smith. Income and wealth inequality have risen in tandem across much of the world since the early 1980s.

\autoref{figure:figure1} and~\autoref{figure:figure2} document both trends using composition-adjusted, GDP-weighted averages across a broad cross-country panel: the top income shares and Gini coefficient have drifted upward over four decades, while public, household, and corporate debt have each expanded substantially. Whether these parallel trends reflect a causal relationship, and if so in which direction and through which channels, remains an open and contested question. In \autoref{table:descriptive}, these trends are tabulated across three periods: the pre-neoliberal era (1970--1978), the immediate post-neoliberal period (1982--1990), and the most recent decade of our estimation window (2010--2022). All series are GDP-weighted composition-adjusted averages constructed to purge sample-entry effects. The periods reflect a structural break in advanced-economy fiscal and financial policy that is well documented in the literature.\footnote{The early 1980s saw a coordinated acceleration of financial liberalization across OECD economies, documented across multiple dimensions in \citet{Abiad2010} and \citet{Quinn1997}. The reintegration of international capital markets following the Bretton Woods era is situated in long-run historical perspective in \citet{ObstfeldTaylor2004}. On the fiscal side, \citet{Alesina1995} show the wave of spending-based consolidation episodes that characterized OECD fiscal policy in this period. The distributional consequences are documented for the United States in \citet{PikettyS2003} and across a broad set of advanced economies in \citet{Atkinson2011}, both of whom identify the early 1980s as the turning point in top income concentration.} Beginning around 1979--1981, advanced economies underwent a well-documented shift toward capital account liberalization, credit and banking deregulation, and a reorientation of fiscal priorities whose effects propagated to the broader cross-country sample through international capital markets. The rise in public debt and the increase in income inequality documented in \autoref{table:descriptive} are both products of this period, and the endogenous relationship between the two has been formalized theoretically by \citet{Azzimonti2014}, who show that rising inequality and financial integration jointly provide incentives to governments to accumulate debt. The break at 1979--1981 accounts for the elections of Margaret Thatcher (1979) and Ronald Reagan (1980) mark a widely recognized structural shift in fiscal and financial policy in advanced economies, whose effects propagated globally through capital market liberalization and the deregulation of credit. Public debt nearly doubled between the pre- and post-neoliberal periods, rising from 24\% to 47\% of GDP on average, while the top 10\% income share rose by nearly 3 percentage points and the top 1\% share by more than 2 points over the same interval. Both debt and inequality continued to rise through the 2010s, with public debt reaching 80\% of GDP and the top 1\% share approaching 17\%.

\begin{table}[ht]
	\centering
	\caption{Debt and Inequality Across Three Periods}
	\label{table:descriptive}
	\begin{threeparttable}
		\begin{tabular}{lccc}
			\toprule
			& Pre-1979 & Post-1981 & Recent \\
			& 1970--1978 & 1982--1990 & 2013--2022 \\
			\midrule
			\addlinespace
			\multicolumn{4}{l}{\textit{Inequality}} \\
			\quad Top 10\% income share      & 36.0\% & 38.9\% & 44.4\% \\
			\quad Top 1\% income share       & 11.3\% & 13.4\% & 16.9\% \\
			\quad Gini coefficient           & 32.5   & 34.3   & 36.0   \\
			\addlinespace
			\multicolumn{4}{l}{\textit{Debt (share of GDP)}} \\
			\quad Public debt                & 24.0\% & 46.6\% & 85.7\% \\
			\quad Household debt             & 40.0\% & 49.5\% & 69.0\% \\
			\quad Non-financial corporate    & 104.9\% & 114.2\% & 139.9\% \\
			\quad Total private debt         & 144.9\% & 163.7\% & 208.9\% \\
			\addlinespace
			\bottomrule
		\end{tabular}
		\begin{tablenotes}
			\scriptsize
			\item \emph{Note}: GDP-weighted averages, composition-adjusted by demeaning within country to remove sample-entry effects. Income shares from WID; Gini from SWIID 9.91; debt from IMF Global Debt Database. Weights are real output-side GDP (PWT 11.0).
		\end{tablenotes}
	\end{threeparttable}
\end{table}

\begin{figure}[ht]
	\caption{Long-Run Trends in Inequality}
	\label{figure:figure1}
	\includegraphics[width=\linewidth]{C:/Users/ageis/Git/chen_twumasi_eisenbarth/figures/inequality_measures.png}
	\caption*{\scriptsize\emph{Note}: GDP-weighted averages, composition-adjusted by demeaning within country. Income shares from WID; Gini from SWIID 9.91. Weights are real output-side GDP (PWT 11.0).}
\end{figure}

The conventional view is that debt can stimulate aggregate demand and output in the short-run, but crowds out capital and reduces output in the long-run (see \cite{Mankiw1999} for a survey). Standard growth models also predict that higher public debt leads to slower growth (see \cite{Saint-Paul91} or \cite{Aizeman07}). There are several channels through which high public debt can adversely affect capital accumulation, productivity and growth: higher long-term interest rates and sovereign risk spillovers to corporate borrowing costs \citep{Gale2003, Cecchetti2015, Corsetti2013}, higher future distortionary taxation and lower future public infrastructure spending, higher inflation \citep{Barro95}, and greater uncertainty about prospects and policies. High debt levels are also likely to constrain the scope for counter-cyclical fiscal policies, which may result in higher volatility and further lower growth \citep{Woo2015}. In more extreme cases of a debt crisis, by triggering a banking or currency crisis, the adverse effects can be magnified \citep{Reinhart2011, Reinhart2012}. The dominant strand of empirical research has framed the fiscal question in terms of aggregate growth, following after \citet{Reinhart2010} and \citet{Reinhart2011}, this  literature has examined whether public debt above some threshold suppresses GDP growth.

\begin{figure}[ht]
	\caption{Long-Run Trends in Debt}
	\label{figure:figure2}
	\includegraphics[width=\linewidth]{C:/Users/ageis/Git/chen_twumasi_eisenbarth/figures/debt_measures.png}
	\caption*{\scriptsize\emph{Note}: GDP-weighted averages, composition-adjusted by demeaning within country. Income shares from WID; Gini from SWIID 9.91. Weights are real output-side GDP (PWT 11.0).}
\end{figure}

The threshold hypothesis has proven fragile under scrutiny. \citet{Eberhardt2015} and \citet{Chudik2017} demonstrate that the debt-growth relationship is characterized by substantial cross-country heterogeneity and that no common threshold exists across economies, even if country-specific tipping points cannot be ruled out. On this point, the empirical relationship between public debt and economic growth in the United States, western offshoot economies, and the OECD countries has also been exhaustively in the past decades (for example \cite{Dwyer82}, \cite{Ahking85}, \cite{Protopapadakis87}, \cite{King1985}, \cite{Woo2015}). These empirical studies provide conflicting results, where some researchers find no significant relation between public debt and growth, while others find some relationship in their analysis with a limited data sample. One possible contributing factor is the paucity of readily available data. Another limitation is the lack of research for developing countries other than the U.S., OECD countries, and their western offshoots.  In other words, an under looked factor is that countries selected for empirical studies in the prior decades are often well-developed countries with strong financial markets and access to foreign financial instruments. Such an arrangement naturally leads to selection bias, limiting the conclusions and inferences that can be drawn from such studies.

The distributional question has attracted far less systematic empirical attention. \citet{Azzimonti2014} show that rising income inequality and financial integration jointly incentivize governments to accumulate debt, suggesting the two trends are endogenously linked rather than causally ordered.  More recently, \citet{Grancini2025} show in a heterogeneous-agent framework that high levels of public debt erode fiscal multipliers primarily through a factor-price channel: elevated real interest rates shift wealth toward households with lower marginal propensities to consume, weakening aggregate demand transmission. On the firm side, \citet{Islam2024} argue that the crowding-out costs of public debt fall disproportionately on small and domestically owned firms, deepening the unequal playing field within the private sector.

These findings share a common implication that motivates this paper: the distributional consequences of public debt are mediated by institutional context, the composition of debt, and the composition of fiscal spending. A country that borrows to finance civilian public investment is likely to generate different distributional outcomes than one that borrows to service existing obligations or fund military expenditure, even at identical debt-to-GDP ratios. Similarly, the consequences of public debt expansion will differ markedly between a country with monetary and fiscal sovereignty operating under a flexible exchange rate and one constrained by fixed rates or external denomination of its obligations. Aggregate debt levels, taken alone, are theoretically insufficient to determine distributional outcomes.

This paper tests that proposition directly. Using an unbalanced panel of 137 countries over the period 1960--2022, we estimate the long-run relationship between public debt and three measures of income inequality, namely the top 10\% income share, the top 1\% income share, and the Gini coefficient, employing the Cross-Sectionally Augmented Distributed Lag (CS-DL) and CS-ARDL estimators of \citet{Chudik2017}, which account for cross-sectional dependence and allow for heterogeneous dynamics across countries. Our baseline results yield null or insignificant effects of public debt on inequality, a finding that can be treated a substantively informative result: debt accumulation in aggregate does not systematically worsen distributional outcomes across a broad cross-country panel. Income group analysis reveals that any detectable debt-inequality relationship is concentrated in lower-middle-income economies, where institutional capacity and fiscal space are most constrained. Fiscal composition analysis, distinguishing civilian government spending from military expenditure and interacting debt levels with spending composition, further shows that what debt finances matters more than its level.


	\section{Distributional Implications of Debt}
	\label{section:section2}
	\subsection{Theoretical Framework}

	An array of research explain debt as the outcome of a political or economic struggle between different groups in the
	population who want to gain more control over resources. The reason debt is accumulated
	is that the group that is in power today may not be in power tomorrow, and debt is a way
	to take advantage of this temporary power. For instance, \cite{Song2012} argue that debt is a tool used to redistribute resources across generations. \cite{Alesina1990} argue that debt represents a way to tie the hands of future governments that will have different preferences from the current one. \cite{Tabellini1991} also illustrates how debt and social security differ as distributional instruments in an overlapping generations environment. In this vein, \cite{Bouton2020} develop a political economic model focusing on the relationship between entitlements and debt. They propose that both debt and entitlements are mechanisms through which temporarily dominant groups can extract resources from groups that will hold power in the future. Specifically, debt reallocates resources across time periods, while entitlements directly influence the future allocation of resources. \cite{Debrun2020} suggest that the presence of endogenous entitlements reduces the incentive for politically powerful groups to accumulate debt, but this leads to an increase in overall government obligations. Furthermore, they argue that fiscal rules can have unintended consequences; if entitlements are not constrained and there are capital market frictions, imposing debt limits can increase total government obligations and result in worse outcomes for all groups.

	A very different approach to understanding public debt is explored by \cite{Azzimonti2014}. They propose a multi-country model with incomplete markets, and they show that governments may choose higher public debt when financial systems are more integrated. \cite{Azzimonti2014} argue that the incentives of governments to borrow rise both as financial markets become internationally integrated and as income inequality increases if this is associated with higher income risk.

	Others argue that public spending financed by government deficits has reduced the inequality gap between low and high-income individuals in their respective countries (see Li and Filer, 2007; Van Bon, 2022). The central argument of this group of research is that "institutional quality" is crucial, and that advanced economies with rule-based governance can effectively utilize public spending to generate net positive outcomes. Countries that have successfully employed fiscal spending have seen it act as a catalyst for social transfers and resource reallocation through economic development (Ortiz-Ospina and Rose, 2016). High public debt, however, can also be seen as a cause of debt crises and political instability. Politicians in high-debt countries may prioritize balancing national accounts due to political pressure from voters, leading to a decrease in public spending policies and a subsequent increase in economic inequality above a certain debt-to-GDP threshold \citep{Reinhart2011}.

	Most existing literature has predominantly examined either the factors influencing debt levels or the effects of overall economic growth on income inequality at the national or global level. However, the lack of studies addressing how debt impacts income inequality limits effective policy making, given the scarcity of knowledge on this direct relationship.

	\subsection{Debt Sustainability and Implications}

	The uncertainty surrounding fiscal policy and the ambiguous contours of a "safe" fiscal space have proven difficult to model, due in large part to the evolving nature of the agents and institutions involved. As \cite{Debrun2019} emphasize, debt sustainability remains an open frontier in macroeconomic analysis, one shaped not only by financial indicators but also by political dynamics. Indeed, sustainability is as much a political question as it is an economic one.

	A central assumption in conventional macroeconomic frameworks is that sustainability is equivalent to solvency; that is, the state’s ability and willingness to meet its future obligations. In this framework, countries can continue to borrow as long as the real interest rate on debt, $r$, is less than the growth rate of the economy, $g$. This condition has formed the basis for a new wave of thinking, most prominently advocated by \cite{Blanchard2019, Blanchard2023} and \cite{FurmanSummers2020}, which suggests that public debt may entail little or no fiscal cost under prevailing conditions. If governments can roll over their debt at low interest rates without raising taxes, the debt-to-GDP ratio may decline over time even in the presence of deficits.

	The empirical literature on the broader social and economic effects of public debt, however, remains contested. Much of the debate centers on whether there exists a specific debt-to-GDP threshold beyond which debt becomes destabilizing. While many studies implicitly or explicitly adopt such a threshold framework, others challenge this view. For instance, \cite{AcalinBall2023} contest the standard postwar narrative that U.S. debt was reduced primarily through economic growth. Instead, they highlight the roles of interest rate pegs, surprise inflation, and tax increases in lowering the debt burden. In their view, the decline in U.S. public debt following World War II was less a function of benign growth dynamics than of active policy interventions.

	On the question of a safe debt level, \cite{Debrun2020} suggest that the United States remains within a "safe zone," even under the possibility of negative shocks. Their estimates indicate that debt levels as high as 160 percent of GDP may be sustainable, so long as real interest rates remain below growth rates and the maturity structure of debt is favorable. \cite{Mian2024} similarly argue that, prior to the COVID-19 pandemic, the U.S. could afford primary deficits above 2 percent of GDP and maintain debt levels over 120 percent of GDP. This view aligns with the broader emerging consensus that advanced economies with deep capital markets and monetary sovereignty face a more flexible debt constraint than previously believed.

	Nonetheless, stabilizing the debt-to-GDP ratio remains a complex endeavor that requires deliberate policy action. The literature highlights several interlocking challenges: how to reconcile debt stabilization with political constraints, how to maintain growth while ensuring the intergenerational fairness of fiscal policy, and how to manage the distributional consequences of tax and spending decisions.

	On one side of the debate, studies such as \cite{Boushey2019}, \cite{Ramey2021}, \cite{Fieldhouse2023}, and \cite{DynanElmendorf2024} emphasize the growth-enhancing effects of public infrastructure investment, particularly when interest rates are low. On the other, researchers such as \cite{YouDutt1996}, \cite{Chetty2017}, \cite{CaseDeaton2023}, and \cite{PikettySaez2019} foreground the distributive implications of fiscal policy, arguing that the composition of spending and taxation plays a pivotal role in shaping both inequality and long-run growth.

	As macroeconomic conditions evolve and policy preferences shift, so too do the frameworks researchers and policymakers use to understand public debt. The literature remains rich and contested, reflecting not only competing economic theories but also the normative judgments that inevitably underlie debates over who pays, who benefits, and what counts as sustainable.




	\section{Related Literature}

	The renowned paper of \cite{Sargent81} discusses
	the "monetary dominance" and "fiscal dominance" regimes said to emerge between the relationship of fiscal deficits  and inflation. Given that the budget deficit is jointly determined by bond sales to the public and seigniorage created by a monetary authority if the monetary authority implements a monetary policy independently, then the fiscal authority faces a budget constraint imposed by the monetary authority when it formulates the fiscal policy. Under this circumstance, the monetary authority can control the money supply, and fiscal deficits do not lead to inflation. In contrast, in a fiscal dominance regime, the monetary authority cannot control the money supply, and fiscal deficits lead to inflation under such fiscal dominance.

	In a widely cited paper, \cite{Reinhart2010} provided long historical data series for the analysis of public-debt-to-GDP ratios and economic growth. Their finding that public-debt-to-GDP ratios above 90\% are associated with markedly lower economic growth rates sparked a major debate. While several leading policymakers in the US and Europe directly referred to \cite{Reinhart2010} in calling for immediate fiscal consolidation measures to reign in public debt, several groups of researchers used the data provided by \cite{Reinhart2010} as well as newly constructed datasets to conduct extensive econometric tests on the impact of public debt levels on economic growth (e.g. \cite{Herndon2014, Woo2015}).  Several papers argue that there is evidence for a negative causal effect of higher public-debt-to-GDP ratios on economic growth (e.g. \cite{Woo2015, Chudik2017}), and for a (close to) 90\% threshold in the public-debt-to-GDP-ratio beyond which growth falls significantly (e.g. \cite{Checherita2012, Baum2013}). While other studies acknowledge the stylized fact of a negative association between initial public debt levels and subsequent growth, they argue that the evidence for a causal effect running from higher public-debt-to-GDP to economic growth is weak at best (e.g. \cite{Panizza2014}). Furthermore, several authors point to systematic differences in the (non-linear) impact of public debt on growth across countries, implying a lack of evidence for universal thresholds in the public-debt-to-GDP ratio beyond which growth falters (e.g. \cite{Eberhardt2015, Eberhardt_2019, Bentour2021}).

	Likewise, the many empirical studies on the relationship between fiscal deficits and inflation provide contradictory and inconsistent results. For the United States, \cite{Hamburger1981} examine the deficit–money relationship from the period of 1954–1976, and conclude that budget deficits are inflationary. The relationship becomes stronger in the "Keynesian period" of 1961--1974. \cite{Dwyer82} uses quarterly data covering the 1953--1978 period to test the relationship between debt, price, and money; his results indicate that expected government deficits have no significance for future inflation. \cite{Ahking85}  examine quarterly data from the period of 1947–1980, and present that the deficit–inflation relationship of the United States does exist during some specific periods. \cite{King1985} investigate the deficit–seigniorage relationship in terms of neoclassical macroeconomic models. They find little connection between fiscal deficits and seigniorage in the 1953–1982 period in the United States; they also estimate the deficit–seigniorage connection of 12 other industrial and developing countries, but still fail to demonstrate that the relationship is significant. \cite{Protopapadakis87} examine the debt–money and the debt–inflation connections for 10 major advanced countries during the period of 1952–1987, and note that the association between debt growth and inflation is very weak. \cite{Haan1990}, investigating 17 developing countries from 1961 to 1985, find no evidence to support the "fiscal dominance hypothesis," discovering that deficits are correlated to inflation during acute inflation periods.

	\cite{Catao2005} model inflation as non-linearly related to fiscal deficits through the inflation tax base and estimate this relationship as intrinsically dynamic, using panel techniques that explicitly distinguish between short- and long-run effects of fiscal deficits. Their results, with a panel spanning 107 countries between 1960 and 2001, give a strong association between deficits and inflation among high-inflation and developing country groups, but not among low-inflation advanced economies.


	In comparison, the relationship between inequality and the level of debt has been studied less. Much of the focus has been on the interplay between household debt and inequality (see for instance, \cite{Kumhof2015, Iacoviello2008, Mason2014}).

\section{Data}
\label{section:data}

The dataset is constructed from six primary sources spanning up to 137 countries over the period 1950--2024, though the effective estimation sample is determined by the minimum consecutive run requirement described below. Country income classifications follow the World Bank Atlas method, assigned using each country's historical modal classification to avoid reclassification artefacts from single-year designations \citep{WorldBankClassifications}. The countries span low, middle, and high income groups across all major regions, providing a broad empirical basis for the analysis. Table~\ref{table:table1} presents the distribution of countries by income group.\footnote{In calculating gross national income (GNI) per capita in U.S. dollars, the World Bank uses the Atlas method. The Atlas conversion factor for a given year is the average of a country's exchange rate for that year and the two preceding years, adjusted for the difference between the rate of inflation in the country and international inflation. GNI in U.S. dollars (Atlas method) for some year $t$ is estimated by dividing the country's GNI in current prices (local currency) by the Atlas factor. The resulting GNI in U.S. dollars can then be divided by a country's midyear population to derive GNI per capita. For 2019, these definitions classify GNI per capita in U.S. dollars of less than 1,035 as “Low Income,” 1,036 to 4,045 as “Lower Middle Income,” 4,046 to 12,535 as “Upper Middle Income,” and GNI per capita greater than 12,535 as “High Income.”}

Because country classifications change over time as GNI per capita evolves, using a single-year designation would introduce discontinuities into the sample stratification: a country that crossed a threshold in a single year would be assigned to different income groups depending on which year was chosen as the reference, with no substantive change in its
structural characteristics. To avoid this, we assign each country its modal classification across all years for which it appears in the historical file. The modal classification represents the income group in which a country spent the largest number of
years over the entire 1987--2024 period, and therefore reflects its dominant structural position rather than its transitory income status in any particular year. Where a country is tied across two groups, the higher group is assigned to avoid under-classifying countries that have experienced sustained development over the sample. The resulting distribution
across the 137 countries in the estimation sample is reported in \autoref{table:table1}.

This approach is preferred to using the most recent classification for two reasons. First, the estimation sample spans several decades, over which many countries have graduated across thresholds; applying the end-of-sample classification would retrospectively impose a later income category on observations from earlier periods when the country had structurally
different characteristics. Second, the research question concerns whether the debt-inequality relationship varies by level of development as a structural feature of the economy, not as a function of where a country happened to sit in a given year relative to a threshold.

\subsection{Real Output and Macroeconomic Aggregates}

Real GDP and related aggregates are drawn from the Penn World Tables (PWT) version 11.0 \citep{FeenstraInklaarTimmer2015}. The primary measure of economic size used in weighting procedures is output-side real GDP (the PWT variable rgdpo), which values production at a country's own final demand prices and is the appropriate concept for cross-country comparisons of productive capacity. GDP per capita and real growth are constructed from this series. The government consumption share of GDP provides the basis for the civilian spending share used in the fiscal decomposition analysis, constructed by subtracting military expenditure from government consumption. Where PWT observations are missing for growth and GDP per capita, we supplement with the corresponding series from the IMF WEO October 2023 vintage \citep{IMFWEO2023}.

\subsection{Public Debt}

The primary source for public debt is the IMF Global Debt Database (GDD) \citep{MbayeMorenoBadiaChae2018}, which provides sectoral debt series for up to 190 countries from 1950, covering general government, central government, household, non-financial corporate, and total private debt, all as a percentage of GDP. We construct the public debt variable using a source hierarchy that prioritizes the most comprehensive institutional coverage: general government debt from the GDD is used where available, supplemented by central government debt from the GDD, and finally by WEO gross general government debt where neither GDD measure is available. This hierarchy follows the approach of \citet{Reinhart2011} while exploiting the broader coverage available in the GDD. We measure public debt as total (domestic plus external) gross central government debt, expressed as a percentage of GDP. For countries where general government debt data is unavailable, we substitute central government debt as a proxy. This approach ensures broader coverage, particularly for developing economies where general government data is often incomplete. The WEO October 2023 vintage \citep{IMFWEO2023} also provides supplementary series for CPI inflation, government expenditure as a share of GDP, and GDP in purchasing power parity terms.

\subsection{Income Inequality}

\subsubsection{Top Income Shares}

The primary inequality measures are top income shares from the World Inequality Database \citep{WID2024}. We use the share of pre-tax national income accruing to the top 10\% and the top 1\% of the national pre-tax income distribution. These series are constructed under the Distributional National Accounts (DINA) framework \citep{BlanchetEtAlDINA2021}, which combines tax records, household surveys, and national accounts to produce distributional estimates that are anchored to macroeconomic aggregates. This anchoring is a crucial methodological advantage: because WID estimates are constrained to sum to national income as measured in the System of National Accounts, they are substantially more comparable across countries and over time than survey-based estimates alone, which are subject to top-coding, non-response at the upper tail, and varying survey methodologies \citep{AtkinsonPiketty2007}. The DINA approach has been applied systematically across more than 100 countries in the World Inequality Report series \citep{ChandelEtAlWIR2022}. For China, WID reports separate series for rural and urban populations; we average these to a single national observation for each year.

Top income shares are theoretically preferred to the Gini coefficient for this paper's research question for two reasons. First, they are more sensitive to changes at the upper tail of the distribution, where the fiscal and financial channels through which public debt operates are most likely to manifest \citep{Atkinson2011, PikettyS2003}. Second, the measurement error in top income shares is better characterized and more tractable than in survey-based Gini estimates: the primary limitation is the degree to which tax records capture top incomes, a concern that is extensively documented and partially addressed in the WID methodology \citep{AlvaradoEtAl2013}.

\subsection*{Gini Coefficient}
The Gini coefficient is drawn from the Standardized World Income Inequality Database, version 9.91 \citep{solt2020}. SWIID provides harmonised estimates of post-tax, post-transfer disposable income inequality for up to 199 countries from 1960 to the present. The harmonization procedure uses the Luxembourg Income Study (LIS) as the comparability benchmark, exploiting its consistent income definition and processing methodology across countries. Observations from other sources (Eurostat, the OECD Income Distribution Database, the World Bank's PovcalNet, the Socio-Economic Database for Latin America and the Caribbean, and national statistical offices) are then expressed relative to their proximate LIS observations using a predictive model that accounts for differences in income definition, unit of observation, and equivalence scale \citep{solt2020}. This approach maximizes comparability without discarding the coverage gains available from non-LIS sources.

Because the harmonization involves imputation across surveys and years, SWIID reports
100 multiple imputations per country-year. We derive a point estimate as the mean Gini
coefficient across all 100 imputations and its standard deviation as a measure of measurement uncertainty. Formally, combining inferences across all 100 imputations following Rubin's rules \citep{Rubin1987} would be the statistically correct treatment; we adopt the simpler averaging procedure for computational tractability and note that the resulting standard errors may modestly understate true uncertainty, particularly in low-income countries where SWIID coverage relies most heavily on imputation across sparse survey years. The Gini coefficient serves as the secondary inequality measure and is used principally in robustness checks; the binding constraint on sample coverage is the WID, which provides wider and more historically continuous country coverage over the relevant sample period.

\subsection{Control Variables}

The choice of control variables follows the theoretical channels identified in the debt and
inequality literatures, and is disciplined by the degrees-of-freedom constraints of the
CS-DL and CS-ARDL estimators.

GDP per capita is included to control for the stage of economic development. The relationship between development and inequality is theoretically ambiguous. The Kuznets hypothesis predicts an inverted-U relationship, while the skill-biased technical change literature predicts monotone increases in the returns to education \citep{KatzMurphy1992, Goldin2008}, but controlling for the level of income is necessary to avoid conflating debt effects with development-stage effects.

Government expenditure as a share of GDP controls for the size of the fiscal state, independent of how it is financed. The direction of the expenditure effect on inequality depends on the composition of spending: social transfers and public health and education
expenditure are associated with reduced inequality \citep{DablaEtAl2015}, while military and general administration spending are not. This motivates the fiscal decomposition
specification, in which civilian government spending (total government expenditure less military spending) is distinguished from the aggregate expenditure share. Military expenditure as a share of GDP is sourced from the World Bank's World Development Indicators \citep{WorldBank2024} and used to construct the civilian spending share.

Life expectancy at birth is included as a broader measure of human development. It captures variation in public health infrastructure, nutrition, and living standards that is not
fully absorbed by GDP per capita or government expenditure, and has been shown to be
associated with income distribution through both human capital accumulation and labor
market participation channels \citep{LondonoPrisca1996}. Life expectancy is preferred to
the PWT human capital index for this purpose because the latter is constructed from assumed Mincer returns to schooling and therefore embeds a theoretical structure that may confound reduced-form estimates of the debt-inequality relationship.

Trade openness, measured as the sum of export and import shares of GDP, is included to control for the distributional consequences of external sector integration. The effects of trade on inequality are theoretically contested: standard Heckscher-Ohlin reasoning predicts that openness reduces inequality in labor-abundant developing countries, while the empirical literature finds mixed results that depend strongly on the skill composition of the labor force and the nature of trading relationships \citep{GoldbergPavcnik2007}.
Excluding trade openness would risk attributing to public debt any distributional effects
that operate through the external sector.

Inflation is included because it is a distributional instrument in its own right. The redistributive incidence of inflation is asymmetric across the wealth distribution: it erodes the real value of nominal assets and liabilities and typically falls more heavily on lower-income households whose wealth is concentrated in liquid deposits rather than real assets \citep{DoepkeSchneider2006}. The composite inflation series draws on WDI CPI inflation, supplemented by WEO CPI where WDI observations are missing. All control variables are sourced from the World Bank's World Development Indicators \citep{WorldBank2024} unless otherwise noted.

\subsection{Exchange Rate Regimes}

Exchange rate regime classifications follow \citet{HarmsKnaze2021}, who construct effective exchange rate regime measures as trade-weighted averages of bilateral regime arrangements. This approach improves on standard unilateral classifications, which assess exchange rate stability only against a single anchor currency, by accounting for a country's full set of trading relationships. The underlying bilateral regime data combine the de-facto classification of \citet{IlzetzkiReinhartRogoff2019} with the IMF's de-jure classification from the Annual Report on Exchange Arrangements and Exchange Restrictions. The resulting measures are available on a coarse 1--4 scale from most rigid to most flexible, covering 1973--2019. We define a floating indicator as 1 if the de-facto regime code corresponds to a managed or free float, a policy autonomy indicator as 1 if the country is not pegged to the US dollar or euro, and a composite monetary sovereignty indicator requiring both conditions to hold in the absence of a contemporaneous currency or debt crisis.

\subsection{Sample Construction}

The estimation sample is determined by identifying, for each country, the longest consecutive run of years with non-missing observations on the core triad: public debt, real GDP growth, and the top 10\% income share. We impose a minimum run length of 30 consecutive years, following \citet{Chudik2017}, who establish this as the minimum time-series dimension required for reliable mean group estimation with the CS-DL and CS-ARDL estimators. Of the 195 countries with any coverage across these sources, 137 satisfy this criterion and form the estimation sample. Their distribution across income groups is reported in \autoref{table:table1}.

Within each country's identified run, isolated missing observations are filled using a centered moving average over a two-period window, applied separately to each variable. This preserves the local trend and autocorrelation properties of each series without imposing a global interpolation model. Countries with gaps too large to bridge under this procedure are excluded from the sample. The excluded set is concentrated among low- and lower-middle-income economies where debt reporting is most limited and includes several high-income economies where reporting discontinuities prevent the construction of a sufficiently long continuous series.


\subsection{Final Dataset}


\begin{table}[htbp]
	\centering
	\caption{Income Classification of Countries}
	\label{table:table1}
	\begin{threeparttable}
		\scriptsize
		\begin{tabular}{llll}
			\toprule
			\textbf{Low Income} & \textbf{Lower Middle Income} & \textbf{Upper Middle Income} & \textbf{High Income} \\
			\midrule
			Bangladesh & Albania & Antigua and Barbuda & Australia \\
			Benin & Algeria & Argentina & Austria \\
			Bhutan & Belize & Azerbaijan & Bahamas \\
			Burkina Faso & Bolivia & Belarus & Bahrain \\
			Burundi & Bulgaria & Botswana & Barbados \\
			Central African Republic & Cabo Verde & Brazil & Belgium \\
			Chad & Colombia & Chile & Canada \\
			Comoros & Congo & Costa Rica & Cyprus \\
			Equatorial Guinea & Cote d'Ivoire & Dominica & Czech Republic \\
			Ethiopia & Ecuador & Gabon & Denmark \\
			Gambia & Egypt & Grenada & Finland \\
			Ghana & El Salvador & Hungary & France \\
			Guinea & Fiji & Kazakhstan & Germany \\
			Guinea-Bissau & Guatemala & Lebanon & Greece \\
			Haiti & Guyana & Malaysia & Iceland \\
			India & Honduras & Mauritius & Ireland \\
			Kenya & Indonesia & Mexico & Israel \\
			Kyrgyzstan & Iran & Oman & Italy \\
			Lao PDR & Jamaica & Panama & Japan \\
			Madagascar & Jordan & Poland & Kuwait \\
			Malawi & Lesotho & Russian Federation & Luxembourg \\
			Mali & Maldives & Saint Kitts and Nevis & Malta \\
			Nepal & Morocco & Saint Lucia & Netherlands \\
			Niger & Namibia & St.\ Vincent and the Grenadines & New Zealand \\
			Nigeria & North Macedonia & Seychelles & Norway \\
			Pakistan & Paraguay & South Africa & Portugal \\
			Rwanda & Philippines & Suriname & Qatar \\
			Sao Tome and Principe & Sri Lanka & Trinidad and Tobago & Saudi Arabia \\
			Senegal & Swaziland & Turkey & Singapore \\
			Sierra Leone & Syrian Arab Republic & Uruguay & Slovakia \\
			Sudan & Thailand & Venezuela & Slovenia \\
			Tanzania & Tunisia & & Spain \\
			Togo & & & Sweden \\
			Uganda & & & Switzerland \\
			Yemen & & & Taiwan \\
			Zambia & & & United Arab Emirates \\
			& & & United Kingdom \\
			& & & United States \\
			\bottomrule
		\end{tabular}
		\begin{tablenotes}
			\scriptsize
			\item \textit{Note:} Income classifications follow the World Bank Atlas method.
			Each country is assigned its modal classification across all years in the World
			Bank Analytical Classifications History file (1987--2021), representing the
			income group in which it spent the largest number of years over the full period
			\citep{WorldBankClassifications}. The sample comprises 137 countries: 36 low
			income, 32 lower middle income, 31 upper middle income, and 38 high income.
		\end{tablenotes}
	\end{threeparttable}
\end{table}

	\section{Methods \& Results}

	\subsection{Debt and Inequality}

	General government debt levels exhibit deep heterogeneity across countries, wealthy and poor, alike. A way of illustrating this is presented in \autoref{figure:figure1}, where we plot general government against against World Bank income classifications defined by \autoref{table:table1}. the distribution of public debt-to-GDP ratios (measuring each country's total government liabilities relative to the size of its economy) exhibits wide internal dispersion within each income bracket. Among low-income countries, debt ratios are generally moderate, although a few outliers suggest emerging fiscal pressures. Lower-middle-income countries display more spread, reflecting diverse macroeconomic conditions: some are undergoing fiscal consolidation, while others may be expanding public borrowing for development purposes or facing debt distress. Upper-middle-income countries, in contrast, show a surprising degree of clustering around relatively lower debt levels, possibly indicative of debt ceilings, external lending constraints, or recent policy reforms.



	\begin{figure}[ht]
		\caption{Gross Government Debt by IMF Income Classification}
		\label{figure:figure3}
		\includegraphics[width=\textwidth]{C:/Users/ageis/Git/chen_twumasi_eisenbarth/figures/total_government_debt_icg.png}
		\caption*{\scriptsize\emph{Note}: Each dot represents a country. Size indicates population. Dotted line is the global average debt ratio; single lines represent income group level debt ratio average. Gross government debt and GDP attained from the IMF WEO from 1980-2019.}
	\end{figure}


	High-income countries present the most striking variation. While some maintain conservative debt profiles, others (including those with advanced welfare states or recent histories of fiscal stimulus) exceed 100 percent of GDP. The wide range among wealthier nations underscores that fiscal capacity and borrowing are not constrained solely by income level but by a complex interplay of historical, institutional, and geopolitical factors.

	Marker size in the figure is proportional to each country’s nominal GDP, allowing for a visual differentiation between economically larger and smaller states within each group. This scaling emphasizes the importance of considering both debt levels and economic size when evaluating fiscal sustainability. Crucially, the figure does not imply that higher income groups uniformly face greater debt burdens, but rather that fiscal heterogeneity is a feature of all income categories, one that defies simplistic generalizations and invites closer examination of national policy regimes, structural constraints, and the political economy of debt accumulation.


	\begin{figure}[ht]
		\caption{Gross Government Debt and Inequality}
		\label{figure:figure4}
		\includegraphics[width=\textwidth]{C:/Users/ageis/Git/chen_twumasi_eisenbarth/figures/figure2.png}
		\caption*{\scriptsize\emph{Note}: Each dot represents a country. Size indicates population. Dotted line is the global average debt ratio; single lines represent income group level debt ratio average. Gross government debt and GDP attained from the IMF WEO from 1980-2019.}
	\end{figure}


	\begin{table}[!h]
		\centering
		\begin{threeparttable}
			\caption{Summary Statistics by Income Group}
			\label{table:table2}
			\fontsize{10}{12}\selectfont
			\begin{tabular}[t]{llrrrrr}
				\toprule
				Income Group & Variable & Observations & Mean & SD & Min & Max \\
				\midrule
				Upper Middle Income & Gini (imputed) & 893 & 46.19 & 6.71 & 26.98 & 63.39 \\
				Upper Middle Income & Government Debt & 893 & 0.41 & 0.28 & 0.04 & 2.38 \\
				Upper Middle Income & GDP & 889 & 0.04 & 0.04 & -0.16 & 0.21 \\
				Upper Middle Income & Population & 893 & 103M & 264M & 657K & 1.41B \\
				High Income & Gini (imputed) & 1410 & 32.12 & 7.02 & 20.52 & 51.53 \\
				High Income & Government Debt & 1410 & 0.56 & 0.36 & 0.04 & 2.36 \\
				High Income & GDP & 1392 & 0.03 & 0.03 & -0.13 & 0.24 \\
				High Income & Population & 1410 & 30.1M & 53.3M & 245K & 328M \\
				Low Income & Gini (imputed) & 470 & 43.86 & 4.13 & 34.76 & 52.16 \\
				Low Income & Government Debt & 470 & 0.50 & 0.31 & 0.03 & 1.69 \\
				Low Income & GDP & 460 & 0.04 & 0.06 & -0.50 & 0.35 \\
				Low Income & Population & 470 & 10.2M & 7.54M & 577K & 42.9M \\
				Lower Middle Income & Gini (imputed) & 1081 & 42.66 & 5.86 & 29.20 & 54.71 \\
				Lower Middle Income & Government Debt & 1081 & 0.52 & 0.32 & 0.02 & 2.61 \\
				Lower Middle Income & GDP & 1077 & 0.04 & 0.04 & -0.22 & 0.26 \\
				Lower Middle Income & Population & 1081 & 77.6M & 205M & 1.14M & 1.38B \\
				Total & Gini (imputed) & 3854 & 39.77 & 8.69 & 20.52 & 63.39 \\
				Total & Government Debt & 3854 & 0.50 & 0.33 & 0.02 & 2.61 \\
				Total & GDP & 3818 & 0.04 & 0.04 & -0.50 & 0.35 \\
				Total & Population & 3854 & 57.9M & 173M & 245K & 1.41B \\
				\bottomrule
			\end{tabular}
			\begin{tablenotes}
				\footnotesize
				\item \textit{Note:} Gini coefficients are from the Standardized World Income Inequality Database (SWIID) v9.3. Government debt and GDP figures are based on IMF World Economic Outlook estimates. Population data are sourced from the World Bank. Observations represent country-year pairs.
			\end{tablenotes}
		\end{threeparttable}
	\end{table}

	\subsection{Estimation}

	In this section we test whether and how central government
	debt as a percentage of GDP changes with income inequality
	in the presence of fixed country and year effects. We first consider the long-run effects of debt and output growth on the distribution of income using a traditional panel fixed-effects approach, in the following equation:
	\begin{equation}
		y_{it} = \alpha_{i} + \lambda_{t} + \gamma d_{it} + \mathbf{\beta}\mathbf{X}_{it} + u_{it}
	\end{equation}
	where $y_{it}$ is the log of the Gini coefficient and $\mathbf{x}_{it}$ a vector of exogenous variables including the the log of debt to GDP ratio $d_{it}$, and $g_{it}$ the growth rate. $\alpha_{i}$ and $\lambda_{t}$ are country and time specific fixed-effects.

	We fit this specification to the entire sample of countries and then separately on each sub-classification of the IMF income classification. These results are presented in \autoref{table:table1}. Column (1) of \autoref{table:table1} is the full dataset, (2) are Upper Middle Income, (3) High Income, (4) Low Income, and (5) Lower Middle Income countries. The largely insignificant findings together with opposite estimated coefficients for two inequality measures challenge the argument by \cite{Azzimonti2014}


	\begin{center}
		\begin{threeparttable}\footnotesize
			\caption{Regression Results by Income Group}
			\label{table:table3}
			\begin{tabular}{@{\extracolsep{5pt}}lccccc}
				\\[-1.8ex]\toprule

				\\[-1.8ex] & (1) & (2) & (3) & (4) & (5)  \\
				\hline \\[-1.8ex]
				Total debt & -0.026 & -0.085 & 0.041 & 0.009 & 0.086$^{***}$ \\
				& (0.067) & (0.059) & (0.041) & (0.029) & (0.021) \\
				& & & & & \\
				GDP & 0.285 & 0.629$^{**}$ & 0.153 & 0.134$^{**}$ & -0.019 \\
				& (0.173) & (0.288) & (0.259) & (0.053) & (0.152) \\
				& & & & & \\
				Human Capital & -0.018 & -0.304 & -0.169$^{**}$ & 0.172$^{***}$ & -0.036 \\
				& (0.137) & (0.184) & (0.066) & (0.051) & (0.099) \\
				& & & & & \\
				Debt $\times$\, GDP & -0.663$^{*}$ & -0.907$^{*}$ & -0.210 & -0.172 & 0.122 \\
				& (0.364) & (0.462) & (0.268) & (0.096) & (0.349) \\
				& & & & &  \\
				\hline \\[-1.8ex]
				Observations & 3,276 & 760 & 1,199 & 397 & 920 \\
				$R^{2}$ & 0.824 & 0.810 & 0.887 & 0.890 & 0.918 \\
				Adjusted $R^{2}$ & 0.818 & 0.794 & 0.880 & 0.875 & 0.912 \\
				\midrule \\
			\end{tabular}
			\noindent
			\begin{tablenotes}
				\item \textit{Note}: All regressions include country and year fixed-effects and are weighted by population. Standard errors clustered at the country level are reported in parentheses. Asterisks denote statistical significance at the 10, 5, and 1 percent levels respectively: {$^{*}$p$<$0.1; $^{**}$p$<$0.05; $^{***}$p$<$0.01}.
			\end{tablenotes}
		\end{threeparttable}
	\end{center}

	\autoref{table:table3} reports the estimates for the full sample and for each income group. Column (1) pools all countries. Growth in domestic product (GDP) enters positively and is marginally significant, indicating that growth is associated with rising inequality on average. The coefficient on debt is close to zero, and human capital has no discernible effect. The interaction term between debt and GDP, however, is negative and weakly significant, suggesting that the inequality-raising effect of growth diminishes when debt levels are higher.

	Column (2), which restricts the sample to upper-middle income countries, shows a clearer pattern. GDP is positively signed and statistically significant at the five-percent level, consistent with the view that growth exacerbates inequality at this stage of development. Debt itself remains insignificant, but the debt–GDP interaction is negative and significant at the ten-percent level, again indicating that debt moderates the distributive impact of growth. Human capital is negatively signed but not significant.

	In column (3), covering high-income countries, the results shift markedly. Neither GDP nor debt displays any independent association with inequality. Human capital, by contrast, is negative and statistically significant, consistent with the interpretation that educational attainment exerts an equalizing influence in advanced economies. The interaction term is small and not significant.

	Column (4), for low-income countries, shows that GDP significantly increases inequality. Unlike in other groups, human capital is positively signed and highly significant, implying that educational investments in these settings may disproportionately benefit elites and thus widen inequality. Debt itself is uncorrelated with inequality, though the debt–GDP interaction is negative and nearly significant, pointing once more to a dampening of growth’s distributive pressures when debt burdens are larger.

	Finally, column (5), which covers lower-middle income economies, demonstrates a distinct pattern: debt is strongly positive and significant at the one-percent level, suggesting a direct inequality-enhancing effect of indebtedness in this group. Neither GDP nor human capital is significant, and the interaction term is uninformative.

	The cross-section of results highlights that the determinants of inequality vary systematically with income level. In the pooled specification (1), and more clearly in the upper-middle income group (2), growth raises inequality, but this effect is attenuated where debt is high. In high-income countries (3), human capital emerges as the main equalizing force, while growth and debt have no role. In low-income economies (4), both growth and human capital increase inequality, with the latter result consistent with theories of elite capture. In lower-middle income countries (5), debt itself is the dominant factor, directly contributing to higher inequality.

	\begin{center}
		\begin{threeparttable}[!htbp]\footnotesize
			\caption{Regression Results}
			\label{table:table4}
			\begin{tabular}{@{\extracolsep{5pt}}lccccc}
				\\[-1.8ex]\toprule

				\\[-1.8ex] & (1) & (2) & (3) & (4) & (5) \\
				\hline \\[-1.8ex]
				Total debt & -0.004 & 0.002 & 0.043$^{**}$ & 0.010 & 0.063$^{**}$  \\
				& (0.045) & (0.081) & (0.025) & (0.031) & (0.031) \\
				& & & & & \\
				GDP & -0.307$^{*}$ & -0.306 & 0.038 & 0.200$^{**}$ & -0.055  \\
				& (0.289) & (0.342) & (0.162) & (0.075) & (0.216) \\
				& & & & & \\
				Human Capital & -0.052 & -0.319$^{**}$ & -0.168$^{***}$ & 0.167$^{**}$ &-0.078 \\
				& (0.130) & (0.017) & (0.152) & (0.062) & (0.101)\\
				& & & & & \\
				Debt $\times$\, GDP & 0.226 & 0.279 & -0.147$^{*}$ & -0.263$^{**}$ & 0.178 \\
				& (0.379) & (0.480) & (0.362) & (0.111) & (0.398) \\
				& & & & &  \\
				\hline \\[-1.8ex]
				Observations & 3,276 & 760 & 1,199 & 397 & 920 \\
				$R^{2}$ & 0.837 & 0.795 & 0.938 & 0.898 & 0.923  \\
				Adjusted $R^{2}$ & 0.831 & 0.777 & 0.934 & 0.882 & 0.917 \\
				\midrule \\
			\end{tabular}
			\noindent
			\begin{tablenotes}
				\item \textit{Note}: All regressions include country and year fixed-effects and weighted by population. Asterisks denote statistical significance at the 10, 5, and 1 percent levels respectively {$^{*}$p$<$0.1; $^{**}$p$<$0.05; $^{***}$p$<$0.01}
			\end{tablenotes}
		\end{threeparttable}
	\end{center}

	The regression estimates provide a nuanced picture. In the aggregate sample, the coefficient on public debt is small and statistically insignificant, suggesting that the level of debt, taken in isolation, does not exert a robust influence on inequality. However, the results diverge once the sample is stratified by income group. Among high-income countries, debt appears positively associated with inequality, whereas in lower-middle-income countries the relationship is also positive but more pronounced. By contrast, in upper-middle-income countries, the estimated association is weak, and in low-income countries the results are imprecise.

	The interaction between debt and output growth provides additional insights. In several specifications, the interaction term carries the opposite sign from the direct effect of debt, implying that the consequences of indebtedness are mediated by the pace of economic expansion. In high-income countries, for example, the interaction is negative, consistent with the view that debt accumulation may exacerbate inequality only in periods of sluggish growth. These heterogeneous results underline the importance of considering both structural conditions and cyclical dynamics when assessing the distributive consequences of public debt.

	\begin{figure}[ht]
		\caption{Gross Government Debt and Inequality}
		\label{figure:figure5}
		\includegraphics[width=\textwidth]{C:/Users/ageis/Git/chen_twumasi_eisenbarth/figures/figure_ME_GDP_debt_panels.png}
		\caption*{\footnotesize\emph{Note}: Panels report the marginal effect of GDP growth on inequality (log Gini) conditional on the debt-to-GDP ratio, with 95 percent confidence intervals (shaded). Darker shading indicates ranges where the confidence interval excludes zero. Dotted vertical lines mark the 10th and 90th percentiles of the observed debt distribution within each income group. Estimates are based on weighted regressions with country and year fixed effects, population weights, and cluster-robust standard errors by country. Data on gross government debt and GDP are from the IMF WEO, 1980–2019.}
	\end{figure}


	The marginal effects plots in \autoref{figure:figure3} clarify the conditional role of public debt in shaping the growth–inequality nexus across the income distribution. In the pooled specification, the slope of the marginal effect curve is clearly downward: at lower levels of indebtedness, increases in GDP are associated with rising inequality, but the effect steadily diminishes as debt rises, eventually turning negative. This pattern is sharpened in the upper-middle income sample, where the inequality-increasing impact of growth at modest debt ratios gives way to a neutral or even inequality-reducing effect as debt burdens mount. The evidence is consistent with the idea that debt acts as a moderator of distributive pressures in transitional economies, mitigating the extent to which growth accrues disproportionately to elites.

	By contrast, the high-income panel shows little evidence of a systematic relationship between growth and inequality across the debt range. The marginal effect curve hovers around zero, with confidence intervals enveloping the axis throughout, suggesting that expansions in advanced economies are distributionally neutral. This echoes the regression estimates, in which human capital rather than growth emerged as the significant equalizing force. In the low-income group, the relationship is more fragile. The marginal effect of growth is again positive at low debt but declines with higher debt, although the confidence intervals are wide. Here the results of the main specification indicate that educational attainment itself is inequality-enhancing, pointing to processes of elite capture that complicate the distributive effects of both growth and schooling in early development. Finally, in the lower-middle income sample the marginal effect of growth is largely flat and insignificant, while debt exerts a direct and substantial inequality-raising influence in the regressions. In these economies, indebtedness appears to drive inequality independently of the growth process, plausibly through the tax and expenditure adjustments that accompany fiscal constraint.

	Taken together, the figures highlight the heterogeneity of mechanisms by stage of development. In the aggregate and especially among upper-middle income countries, growth is inequality-increasing at low debt levels, but this effect is attenuated—and in some cases reversed—when debt is high. In rich economies, growth plays no role, with educational institutions instead determining the distributional trajectory. In poor countries, both growth and human capital risk amplifying inequality under elite control. And among lower-middle income economies, debt itself emerges as the critical determinant of inequality. The comparative plots thus reinforce the regression results while offering a visual demonstration of how debt conditions the growth–inequality relationship in markedly different ways across the income spectrum.

	\subsection{Dynamic Analysis: ARDL Bounds Testing Approach}

	While the fixed-effects estimates provide useful baseline associations between debt, growth, and inequality, they assume a static long-run relationship and do not account for the dynamic adjustment processes that characterize macroeconomic variables. To address these limitations, we employ the autoregressive distributed lag (ARDL) bounds testing approach developed by \cite{PesaranShin1999} and \cite{PesaranShinSmith2001}. This methodology offers several advantages for our analysis: it permits the inclusion of variables with different orders of integration (I(0) or I(1)), it performs well in small samples, and it allows for the estimation of both short-run dynamics and long-run equilibrium relationships.

	When estimating the distributional consequences of sovereign debt and economic growth, addressing simultaneous causality and endogenous feedback loops is paramount. Traditional instrumental variable approaches (FE-2SLS) require strictly exogenous external instruments, a condition that is practically elusive in a diverse, 82-country macroeconomic panel, as global shocks rarely satisfy the exclusion restriction across universally heterogeneous economies. Similarly, while system GMM is frequently employed to resolve dynamic panel endogeneity, it is mathematically optimized for ``large $N$, small $T$'' datasets. Applying system GMM to a macroeconomic panel with a sufficiently long time dimension risks severe instrument proliferation, which overfits the endogenous variables and artificially inflates Hansen $J$-test $p$-values. To resolve these dual challenges without relying on arbitrary external instruments or risking GMM instrument proliferation, this study employs the panel ARDL framework. As demonstrated by \cite{PesaranShin1999}, augmenting the ARDL specification with an appropriate lag structure for both the dependent and explanatory variables fully absorbs residual serial correlation. Crucially, provided the variables exhibit a valid cointegrating relationship, the ARDL framework inherently corrects for endogenous regressors, yielding ``super-consistent'' long-run parameter estimates that are asymptotically immune to simultaneous causality bias.

	The ARDL model for each country takes the general form:
	\begin{equation}
		\Delta y_{t} = \alpha + \sum_{j=1}^{p} \phi_{j} \Delta y_{t-j} + \sum_{j=0}^{q_{1}} \beta_{1,j} \Delta d_{t-j} + \sum_{j=0}^{q_{2}} \beta_{2,j} \Delta g_{t-j} + \sum_{j=0}^{q_{3}} \beta_{3,j} \Delta h_{t-j} + \sum_{j=0}^{q_{4}} \beta_{4,j} \Delta (d \times g)_{t-j} + \lambda \text{ECT}_{t-1} + \varepsilon_{t}
	\end{equation}

	where $y_{t}$ denotes the log Gini coefficient, $d_{t}$ is the debt-to-GDP ratio, $g_{t}$ is GDP growth, $h_{t}$ is human capital, and $(d \times g)_{t}$ captures the interaction between debt and growth. The error correction term $\text{ECT}_{t-1}$ represents the speed of adjustment toward the long-run equilibrium, and $\lambda$ is expected to be negative and significant when a cointegrating relationship exists.

	We estimate country-specific ARDL models, selecting the optimal lag structure using the Akaike Information Criterion (AIC) with a maximum lag order of five. This approach respects the heterogeneity emphasized throughout our analysis while providing estimates of both short-run responses and long-run multipliers for each country.

	\subsubsection{Cointegration Evidence}

	The ARDL bounds test evaluates the null hypothesis of no long-run relationship against the alternative of cointegration. \autoref{table:bounds_summary} summarizes the distribution of bounds test results across the 81 countries in our sample.

	\begin{table}[htbp]
		\centering
		\begin{threeparttable}
			\caption{Summary of ARDL Bounds Test Results}
			\label{table:bounds_summary}
			\begin{small}
				\begin{tabular}{lcc}
					\toprule
					\textbf{Conclusion} & \textbf{Number of Countries} & \textbf{Percent} \\
					\midrule
					Strong evidence of cointegration ($p \leq 0.01$) & 25 & 30.9 \\
					Evidence of cointegration ($p \leq 0.05$) & 10 & 12.3 \\
					Weak evidence of cointegration ($0.05 < p \leq 0.10$) & 5 & 6.2 \\
					No evidence of cointegration ($p > 0.10$) & 31 & 38.3 \\
					Inconclusive / Error & 10 & 12.3 \\
					\bottomrule
				\end{tabular}
			\end{small}
			\begin{tablenotes}
				\footnotesize
				\item \textit{Note}: Results based on Pesaran-Shin-Smith bounds tests. Critical values from \citet{PesaranShinSmith2001} assuming unrestricted intercept and no trend. Countries with computational errors are excluded from the cointegration subsample analysis.
			\end{tablenotes}
		\end{threeparttable}
	\end{table}

	The bounds tests reveal that approximately half of the countries in our sample (40 out of 71 with valid tests, or 56 percent) exhibit at least weak evidence of a long-run equilibrium relationship between inequality, debt, growth, and human capital. This finding supports the use of an error correction framework for these countries, as the variables do not drift apart indefinitely but instead adjust toward a stable long-run relationship.

	\autoref{table:bounds_countries} presents selected country-level bounds test results, organized by income classification. Among high-income countries, strong evidence of cointegration is found in Australia, Austria, Belgium, France, Luxembourg, South Korea, Spain, Sweden, and Switzerland. In the middle-income groups, China, Colombia, Costa Rica, Ecuador, Egypt, and South Africa display robust cointegrating relationships. Among lower-income economies, Bangladesh, Ghana, Kenya, Malawi, and Madagascar also show strong evidence of long-run equilibrium.

	\begin{table}[htbp]
		\centering
		\begin{threeparttable}
			\caption{Bounds Test Results by Income Group (Selected Countries)}
			\label{table:bounds_countries}
			\begin{small}
				\begin{tabular}{llcc}
					\toprule
					\textbf{Income Group} & \textbf{Country} & \textbf{F-Statistic} & \textbf{Conclusion} \\
					\midrule
					\multirow{5}{*}{High Income}
					& Belgium & 9.60 & Strong ($p \leq 0.01$) \\
					& Luxembourg & 10.19 & Strong ($p \leq 0.01$) \\
					& Sweden & 7.71 & Strong ($p \leq 0.01$) \\
					& United States & 4.94 & Evidence ($p \leq 0.05$) \\
					& United Kingdom & 4.44 & Evidence ($p \leq 0.05$) \\
					\midrule
					\multirow{5}{*}{Upper Middle Income}
					& China & 8.91 & Strong ($p \leq 0.01$) \\
					& South Africa & 7.71 & Strong ($p \leq 0.01$) \\
					& Costa Rica & 5.80 & Strong ($p \leq 0.01$) \\
					& Ecuador & 6.35 & Strong ($p \leq 0.01$) \\
					& Brazil & 3.92 & Weak ($0.05 < p \leq 0.10$) \\
					\midrule
					\multirow{5}{*}{Lower Middle Income}
					& Ghana & 10.51 & Strong ($p \leq 0.01$) \\
					& Egypt & 7.07 & Strong ($p \leq 0.01$) \\
					& Philippines & 7.89 & Strong ($p \leq 0.01$) \\
					& Bangladesh & 5.83 & Strong ($p \leq 0.01$) \\
					& Kenya & 5.69 & Strong ($p \leq 0.01$) \\
					\midrule
					\multirow{3}{*}{Low Income}
					& Madagascar & 9.40 & Strong ($p \leq 0.01$) \\
					& Malawi & 8.40 & Strong ($p \leq 0.01$) \\
					& Uganda & 4.47 & Evidence ($p \leq 0.05$) \\
					\bottomrule
				\end{tabular}
			\end{small}
			\begin{tablenotes}
				\footnotesize
				\item \textit{Note}: F-statistics from ARDL bounds tests. Critical values at 1\%, 5\%, and 10\% significance levels are approximately 4.29, 3.23, and 2.75 respectively for the case with unrestricted intercept and no trend (k=4 regressors).
			\end{tablenotes}
		\end{threeparttable}
	\end{table}

	The cross-sectional heterogeneity in cointegration evidence is notable. Countries where no cointegrating relationship is detected include several with historically volatile economic conditions (Argentina, Bolivia, Jamaica, Zimbabwe) as well as some stable advanced economies (Canada, Japan, Netherlands, Norway). This pattern suggests that the existence of a long-run equilibrium is not simply a function of development level but depends on country-specific institutional and structural factors.

	\subsubsection{Long-Run Multipliers}

	For countries where cointegration is established, the long-run multipliers capture the eventual cumulative impact of a permanent change in each explanatory variable on inequality. \autoref{table:meta_longrun} presents the meta-analytic summary of these multipliers across the 81 country-specific ARDL models.

	\begin{table}[htbp]
		\centering
		\begin{threeparttable}
			\caption{Meta-Analysis of Long-Run Multipliers}
			\label{table:meta_longrun}
			\begin{small}
				\begin{tabular}{lcccccc}
					\toprule
					\textbf{Variable} & \textbf{Mean} & \textbf{SD} & \textbf{Median} & \textbf{Min} & \textbf{Max} & \textbf{N} \\
					\midrule
					Debt-to-GDP & $-0.41$ & 6.14 & 0.006 & $-31.18$ & 34.63 & 81 \\
					GDP Growth & 9.50 & 124.49 & 0.010 & $-134.11$ & 1101.62 & 81 \\
					Human Capital & 0.75 & 3.14 & 0.002 & $-1.97$ & 15.02 & 81 \\
					Debt $\times$ Growth & $-6.93$ & 106.65 & 0.247 & $-896.35$ & 284.83 & 81 \\
					\bottomrule
				\end{tabular}
			\end{small}
			\begin{tablenotes}
				\footnotesize
				\item \textit{Note}: Long-run multipliers computed as $-\sum \beta_j / (1 - \sum \phi_j)$ from country-specific ARDL models. The dependent variable is the log Gini coefficient. Large standard deviations reflect substantial cross-country heterogeneity.
			\end{tablenotes}
		\end{threeparttable}
	\end{table}

	The most striking feature of \autoref{table:meta_longrun} is the substantial heterogeneity in long-run effects across countries. For all four variables, the standard deviations vastly exceed the means, and the ranges span from large negative to large positive values. This dispersion underscores the limitations of any single ``global'' estimate of the debt-inequality relationship and reinforces the importance of country-specific analysis.

	Focusing on the median values, which are less sensitive to outliers, we find that the typical long-run effect of debt on inequality is essentially zero (0.006), as is the typical effect of GDP growth (0.010). Human capital similarly shows a near-zero median effect (0.002), while the debt-growth interaction has a small positive median (0.247). These results suggest that, for the median country, there is no systematic long-run relationship between debt accumulation and income inequality, a finding consistent with the insignificant pooled coefficients in our static fixed-effects specifications.

	However, the extreme values in the distribution warrant attention. In some countries (such as the Philippines and Belgium), debt appears to have substantial long-run effects on inequality, though with opposite signs. The Belgian case, where the debt multiplier reaches 34.6, reflects a combination of high short-run responsiveness and slow adjustment dynamics, magnifying the cumulative effect. Conversely, the large negative multiplier for the Philippines ($-31.2$) suggests that debt accumulation is associated with declining inequality over the long run in that context, potentially through channels of public investment or social spending.

	\subsubsection{Error Correction Dynamics}

	The error correction term (ECT) coefficient measures the speed at which the system returns to equilibrium following a shock. A negative and significant ECT indicates stable adjustment dynamics, while a positive coefficient suggests explosive behavior incompatible with a stable long-run relationship.

	\autoref{table:ecm_summary} summarizes the distribution of ECT coefficients across countries.

	\begin{table}[htbp]
		\centering
		\begin{threeparttable}
			\caption{Distribution of Error Correction Coefficients}
			\label{table:ecm_summary}
			\begin{small}
				\begin{tabular}{lcc}
					\toprule
					\textbf{ECT Range} & \textbf{Number of Countries} & \textbf{Interpretation} \\
					\midrule
					$\lambda < -0.30$ & 9 & Fast adjustment (< 3 years) \\
					$-0.30 \leq \lambda < -0.10$ & 24 & Moderate adjustment (3--10 years) \\
					$-0.10 \leq \lambda < 0$ & 34 & Slow adjustment (> 10 years) \\
					$\lambda \geq 0$ & 14 & No stable adjustment \\
					\midrule
					\textbf{Total} & \textbf{81} & \\
					\bottomrule
				\end{tabular}
			\end{small}
			\begin{tablenotes}
				\footnotesize
				\item \textit{Note}: The half-life of adjustment is approximated by $\ln(0.5)/\ln(1+\lambda)$. Positive ECT values indicate absence of error correction dynamics.
			\end{tablenotes}
		\end{threeparttable}
	\end{table}

	The majority of countries (67 of 81) display negative ECT coefficients, consistent with stable long-run dynamics. Among these, the distribution is skewed toward slow adjustment: 34 countries have ECT values between $-0.10$ and zero, implying half-lives of adjustment exceeding a decade. This sluggish pace of equilibrium restoration is economically intuitive: income distributions are shaped by deep structural factors (labor market institutions, educational systems, tax policies) that change only gradually.

	Countries with faster adjustment speeds (ECT $< -0.30$) include Egypt ($-0.42$), Malaysia ($-0.41$), Morocco ($-0.42$), Sweden ($-0.55$), and France ($-0.99$). The French case is particularly notable: an ECT of $-0.99$ implies near-complete adjustment within a single year, suggesting either highly responsive policy mechanisms or measurement idiosyncrasies. At the other extreme, countries such as Argentina, Algeria, and the Philippines show ECT values near zero or slightly positive, indicating either no error correction or very weak mean-reversion.

	\autoref{table:ecm_selected} presents ECT coefficients and model fit statistics for selected countries.

	\begin{table}[htbp]
		\centering
		\begin{threeparttable}
			\caption{Error Correction Coefficients: Selected Countries}
			\label{table:ecm_selected}
			\begin{small}
				\begin{tabular}{llccc}
					\toprule
					\textbf{Income Group} & \textbf{Country} & \textbf{ECT ($\lambda$)} & \textbf{ECM $R^2$} & \textbf{Half-Life (years)} \\
					\midrule
					\multirow{5}{*}{High Income}
					& France & -0.99 & 0.87 & 0.2 \\
					& Sweden & -0.55 & 0.83 & 0.9 \\
					& Spain & -0.28 & 0.59 & 2.1 \\
					& South Korea & -0.34 & 0.72 & 1.7 \\
					& United States & -0.07 & 0.55 & 9.5 \\
					\midrule
					\multirow{4}{*}{Upper Middle Income}
					& China & -0.23 & 0.97 & 2.6 \\
					& South Africa & 0.05 & 0.65 & --- \\
					& Brazil & -0.07 & 0.64 & 9.6 \\
					& Costa Rica & -0.13 & 0.79 & 4.9 \\
					\midrule
					\multirow{4}{*}{Lower Middle Income}
					& Egypt & -0.42 & 0.65 & 1.3 \\
					& Malaysia & -0.41 & 0.62 & 1.3 \\
					& India & -0.12 & 0.33 & 5.4 \\
					& Ghana & -0.14 & 0.63 & 4.6 \\
					\midrule
					\multirow{3}{*}{Low Income}
					& Madagascar & -0.26 & 0.82 & 2.3 \\
					& Malawi & -0.23 & 0.91 & 2.7 \\
					& Mali & -0.21 & 0.56 & 2.9 \\
					\bottomrule
				\end{tabular}
			\end{small}
			\begin{tablenotes}
				\footnotesize
				\item \textit{Note}: Half-life computed as $\ln(0.5)/\ln(1+\lambda)$. ``---'' indicates no stable equilibrium (positive ECT). ECM $R^2$ is the variance explained by the error correction model.

			\end{tablenotes}
		\end{threeparttable}
	\end{table}

	The ECM $R^2$ values provide an additional diagnostic. Countries with high explanatory power (China at 0.97, Malawi at 0.91, Argentina at 0.99) suggest that the ARDL specification captures the essential dynamics of inequality in those contexts. Lower values (India at 0.33, Mauritius at 0.07) indicate that additional factors beyond debt, growth, and human capital drive inequality dynamics in those settings.

	\subsection{Synthesis and Robustness}

	The dynamic ARDL analysis complements and extends the static fixed-effects results in several important respects. First, it confirms that no universal relationship exists between public debt and income inequality. The pooled and median estimates are indistinguishable from zero, and the enormous cross-country variance in long-run multipliers defies any single characterization.

	Second, the cointegration tests provide formal evidence that, for roughly half of the countries studied, debt, growth, human capital, and inequality share a long-run equilibrium. This finding supports the theoretical premise that fiscal variables and distributional outcomes are jointly determined in the long run, even if short-run correlations are weak or absent.

	Third, the error correction analysis reveals that adjustment toward equilibrium is generally slow. Where stable dynamics exist, the typical half-life exceeds a decade, implying that the distributional consequences of fiscal policy unfold over very long horizons. This has direct policy implications: short-run evaluations of debt-financed programs may miss their full redistributive effects.

	Finally, the substantial heterogeneity documented throughout this analysis reinforces our central argument against universal debt thresholds. The relationship between debt and inequality varies not only in magnitude but in sign across countries, and these differences are systematic with respect to income level, institutional context, and exchange rate regime. A policy framework that ignores this heterogeneity risks prescribing remedies ill-suited to local conditions.

	\subsection{Exchange Rate Regimes and Monetary Sovereignty}

	A central motivation for this study is the hypothesis that countries with monetary sovereignty, defined as possessing an independent monetary policy and a floating exchange rate, may experience fundamentally different debt-inequality dynamics than those operating under fixed or pegged arrangements. This conjecture draws on insights from both heterodox monetary theory and the political economy of fiscal policy. If monetarily sovereign nations face softer budget constraints and can deploy deficit spending without the disciplining pressures of currency crises or external creditors, they may use fiscal space more expansively and, potentially, more progressively.

	To test this hypothesis, we draw on exchange rate regime classifications from \citet{HarmsKnaze2021}, which extend the \citet{IlzetzkiReinhartRogoff2019} methodology to construct both de jure and de facto measures of regime flexibility. We define three key indicators: (1) \textit{floating}, coded as 1 if a country's de facto IRR regime code is 12 or higher (managed or free float); (2) \textit{policy autonomy}, coded as 1 if the country is not pegged to the U.S. dollar or euro; and (3) \textit{meets hypothesis}, coded as 1 if a country satisfies both conditions and is not experiencing a currency or debt crisis. We also employ a continuous measure of effective regime flexibility on a 1--4 scale, capturing trade-weighted exchange rate arrangements.

	\subsubsection{Panel Evidence with Regime Interactions}

	We augment our baseline fixed-effects specification with interaction terms between debt and the regime indicators. \autoref{table:regime_interactions} presents the results.

	\begin{table}[htbp]
		\centering
		\begin{threeparttable}
			\caption{Debt-Inequality Relationship by Exchange Rate Regime}
			\label{table:regime_interactions}
			\begin{small}
				\begin{tabular}{lcccccc}
					\toprule
					& (1) & (2) & (3) & (4) & (5) & (6) \\
					& Base & Floating & Eff. Regime & Autonomy & Hypothesis & Triple \\
					\midrule
					Debt & 0.033 & 0.017 & $-0.077$ & 0.026 & 0.022 & 0.018 \\
					& (0.022) & (0.018) & (0.056) & (0.018) & (0.017) & (0.017) \\
					GDP & $-0.044$ & $-0.006$ & $-0.036$ & $-0.047$ & $-0.060$ & $-0.087$ \\
					& (0.071) & (0.066) & (0.065) & (0.067) & (0.065) & (0.063) \\
					Human Capital & $-0.034$ & $-0.032$ & $-0.030$ & $-0.037$ & $-0.032$ & $-0.034$ \\
					& (0.044) & (0.044) & (0.043) & (0.045) & (0.043) & (0.043) \\
					Debt $\times$ GDP & 0.018 & 0.089 & 0.006 & 0.020 & 0.056 & \\
					& (0.109) & (0.092) & (0.092) & (0.099) & (0.093) & \\
					Debt $\times$ Floating & & 0.045$^{**}$ & & & & \\
					& & (0.021) & & & & \\
					Debt $\times$ Eff. Regime & & & 0.037$^{**}$ & & & \\
					& & & (0.018) & & & \\
					Debt $\times$ Autonomy & & & & 0.023 & & \\
					& & & & (0.037) & & \\
					Debt $\times$ Hypothesis & & & & & 0.048 & \\
					& & & & & (0.032) & \\
					Debt $\times$ Meets Hyp. & & & & & & 0.083$^{**}$ \\
					& & & & & & (0.036) \\
					\midrule
					Country FE & Yes & Yes & Yes & Yes & Yes & Yes \\
					Year FE & Yes & Yes & Yes & Yes & Yes & Yes \\
					Observations & 3,209 & 3,209 & 3,209 & 3,209 & 3,209 & 3,209 \\
					$R^2$ & 0.948 & 0.950 & 0.950 & 0.949 & 0.949 & 0.949 \\
					Within $R^2$ & 0.025 & 0.050 & 0.046 & 0.031 & 0.038 & 0.043 \\
					\bottomrule
				\end{tabular}
			\end{small}
			\begin{tablenotes}
				\footnotesize
				\item \textit{Note}: Dependent variable is log Gini coefficient. Standard errors clustered by country in parentheses. ``Floating'' indicates IRR regime $\geq 12$; ``Eff. Regime'' is continuous 1--4 scale; ``Autonomy'' indicates no USD/EUR peg; ``Meets Hyp.'' requires floating, autonomy, and no crisis. $^{*}p<0.1$; $^{**}p<0.05$; $^{***}p<0.01$.
			\end{tablenotes}
		\end{threeparttable}
	\end{table}

	The results in \autoref{table:regime_interactions} reveal a pattern that challenges the initial hypothesis. The interaction between debt and floating exchange rates (column 2) is positive and statistically significant at the five-percent level ($\beta = 0.045$, $p = 0.032$), indicating that countries with flexible exchange rates experience a \textit{stronger} positive relationship between debt and inequality than those with fixed arrangements. The same pattern emerges with the continuous effective regime measure (column 3): more flexible regimes are associated with larger inequality-increasing effects of debt accumulation ($\beta = 0.037$, $p < 0.05$).

	The full hypothesis test in column 6, which interacts debt with an indicator for countries satisfying all sovereignty conditions, confirms this finding. The coefficient on Debt $\times$ Meets Hypothesis is 0.083 and significant at the five-percent level, implying that monetarily sovereign countries exhibit debt-inequality dynamics approximately 0.08 log points more positive than constrained countries—the opposite of what the sovereignty hypothesis predicts.

	\subsubsection{Split-Sample Analysis}

	To probe these findings further, we estimate separate models for floating and fixed-regime countries. The split-sample coefficients on debt are:

	\begin{center}
		\begin{tabular}{lcc}
			\toprule
			\textbf{Sample} & \textbf{Debt Coefficient} & \textbf{N} \\
			\midrule
			Floating regimes & 0.033 & 1,247 \\
			Fixed regimes & 0.020 & 1,962 \\
			Meets hypothesis & 0.032 & 892 \\
			Fails hypothesis & 0.024 & 2,317 \\
			\bottomrule
		\end{tabular}
	\end{center}

	In all cases, the debt coefficient is larger (more inequality-increasing) for sovereign countries than for constrained ones, though the differences are modest in magnitude and the individual coefficients are not statistically distinguishable from zero.

	\subsubsection{Long-Run Dynamics by Regime}

	The short-run panel results, however, may not capture the full story. If monetary sovereignty enables policies whose redistributive benefits materialize only gradually—through sustained public investment, education spending, or social insurance—then the long-run multipliers from our ARDL models may reveal a different pattern.

	\autoref{table:lr_regime} compares the distribution of long-run debt multipliers across regime categories.

	\begin{table}[htbp]
		\centering
		\begin{threeparttable}
			\caption{Long-Run Debt Multipliers by Exchange Rate Regime}
			\label{table:lr_regime}
			\begin{small}
				\begin{tabular}{lccccc}
					\toprule
					\textbf{Regime} & \textbf{N} & \textbf{Mean LR} & \textbf{Difference} & \textbf{$t$-stat} & \textbf{$p$-value} \\
					\midrule
					Floating (IRR $\geq$ 12) & 12 & $-0.302$ & & & \\
					Fixed (IRR $<$ 12) & 53 & $-0.005$ & $-0.297$ & $-0.82$ & 0.418 \\
					\midrule
					Meets Hypothesis & 5 & $-0.150$ & & & \\
					Fails Hypothesis & 60 & $-0.053$ & $-0.097$ & $-0.18$ & 0.860 \\
					\bottomrule
				\end{tabular}
			\end{small}
			\begin{tablenotes}
				\footnotesize
				\item \textit{Note}: Long-run multipliers from country-specific ARDL models. Floating defined as IRR regime code $\geq 12$. ``Meets Hypothesis'' requires floating regime, policy autonomy, and no contemporaneous crisis. $t$-statistics from two-sample $t$-tests assuming unequal variances.
			\end{tablenotes}
		\end{threeparttable}
	\end{table}

	The long-run estimates present a strikingly different picture. Among floating-regime countries, the mean long-run debt multiplier is $-0.302$, compared to $-0.005$ for fixed-regime countries. Countries meeting the full sovereignty hypothesis show a mean multiplier of $-0.150$ versus $-0.053$ for those that do not. In both cases, the direction of the difference \textit{reverses}: sovereign countries appear to exhibit more inequality-\textit{reducing} long-run effects of debt, not more inequality-increasing effects as the short-run panel estimates suggest.

	However, neither difference is statistically significant. The $t$-statistics are well below conventional thresholds ($-0.82$ and $-0.18$), and the $p$-values (0.418 and 0.860) provide no grounds for rejecting the null of equal means. The failure to detect significance reflects two factors: the small number of countries classified as floating or sovereignty-compliant (12 and 5, respectively), and the enormous cross-country variance in long-run multipliers documented earlier.

	\subsubsection{Interpretation}

	The regime analysis yields findings that resist easy characterization. In the short run, the evidence favors the mainstream view: countries with monetary sovereignty and floating exchange rates experience \textit{stronger} positive associations between debt and inequality. This pattern is consistent with several mechanisms. Sovereign countries may deploy debt in ways that disproportionately benefit asset holders—through financial sector support, quantitative easing, or tax policies that accompany deficit spending. The absence of external discipline does not, by itself, translate into progressive fiscal choices.

	In the long run, the pattern reverses direction, though not significantly. Sovereign countries show more negative (inequality-reducing) debt multipliers on average, hinting that the full redistributive consequences of fiscal policy may take decades to materialize. But the small samples and wide confidence intervals preclude strong conclusions.

	From a theoretical standpoint, these results offer partial support for both camps. The mainstream concern that monetary sovereignty does not guarantee equitable outcomes finds confirmation in the short-run dynamics. The heterodox intuition that sovereign countries \textit{could} use fiscal space progressively is not refuted by the long-run estimates—but neither is it affirmed. The critical insight may be that sovereignty creates \textit{capacity} without determining \textit{outcomes}: what matters is not whether a country can run deficits freely, but how it chooses to spend.


	\subsection{Discussion}

	Taken together, the results complicate simple narratives that treat public debt as uniformly detrimental or uniformly benign. The relationship is conditional: it depends on the stage of development, institutional capacity, and the growth environment in which debt is accumulated. For advanced economies with well-developed financial markets, rising debt levels may coincide with higher inequality, potentially reflecting the regressive incidence of debt-financed fiscal consolidation, or the unequal distributional effects of accommodative monetary policies that accompany high debt ratios. For lower-middle-income countries, debt accumulation may exacerbate inequality through more direct channels, including increased exposure to external shocks and the need for austerity measures imposed by international creditors.

	By contrast, the absence of a strong relationship in upper-middle-income countries suggests that institutional factors (such as targeted social transfers, fiscal rules, or debt ceilings) may mitigate the distributive consequences of public borrowing. The evidence that human capital investments reduce inequality in wealthier states but exacerbate disparities in poorer ones is particularly instructive. It points to the role of access: where educational and health resources are broadly available, public debt that finances these expenditures has equalizing effects; where access is restricted, debt-financed investments can reinforce existing divides.

	Our findings also speak to the broader debate on fiscal space and debt sustainability. If the distributive consequences of debt vary so markedly across income groups, then one-size-fits-all policy prescriptions, such as universal debt thresholds, are misleading. Instead, the political economy context (monetary sovereignty, exchange rate regime, and institutional quality) determines whether public debt functions as a stabilizing or destabilizing force in the long run.


	\section{Conclusion}

	This paper has examined the complex interplay between public debt, economic growth, and income inequality across 82 countries over nearly five decades. Contrary to the notion of a universal debt threshold beyond which growth and equity outcomes deteriorate, our results highlight deep heterogeneity. Debt does not exert a uniform effect on inequality; instead, its distributive impact depends on structural conditions, the level of development, the interaction with economic growth, and the institutional environment in which fiscal policy operates.

	The central finding of this study is that the transmission of debt into inequality is conditional on institutional quality. This leads us to advance an institutional-first approach to debt management: public debt should not be treated as a universal stabilization tool, but as an instrument that requires a specific level of administrative and institutional maturity to operate safely. The heterogeneity documented throughout our analysis is not a statistical nuisance; it is the result itself. Where institutions are robust (where tax administration is effective, budget processes are transparent, and central banks operate with credible independence) debt can serve as a macro-stabilizer during downturns, channeling resources toward social protection and productive investment. Where these foundations are absent, the same fiscal instrument risks amplifying inequality through elite capture, regressive consolidation, and external vulnerability.

	Three specific conclusions follow. First, debt accumulation is most likely to exacerbate inequality in lower-middle-income and high-income economies, albeit through different mechanisms. In lower-middle-income countries, the "austerity channel" dominates: weak institutional capacity means that debt obligations translate into expenditure cuts that fall disproportionately on the poor, while external creditor pressures further constrain progressive fiscal responses. In high-income economies, the inequality-increasing effects of debt operate through more subtle channels, regressive fiscal adjustments and the unequal distributional consequences of accommodative monetary policies that accompany high debt ratios. For both groups, the implication is the same: institutional context determines whether debt is stabilizing or destabilizing.

	Second, economic growth on its own is not necessarily equalizing. In low-income countries, growth often coincides with widening disparities, and our finding that human capital investments can exacerbate inequality in these settings points to processes of elite capture. Where access to education and health resources remains stratified, debt-financed investments reinforce existing divides rather than narrowing them. This underscores the need for institutional prerequisites (broad-based access to public services, progressive tax structures, and accountable governance) before debt can function as a vehicle for inclusive development.

	Third, the exchange rate regime analysis reveals that monetary sovereignty creates fiscal \textit{capacity} without determining fiscal \textit{outcomes}. In the short run, sovereign countries with floating exchange rates actually exhibit stronger positive associations between debt and inequality. The long-run ARDL estimates hint at a reversal, with sovereign countries showing more inequality-reducing debt multipliers on average, but the evidence is not statistically decisive. The critical insight is that the absence of external discipline does not, by itself, translate into progressive fiscal choices. What matters is not whether a country \textit{can} run deficits freely, but how it \textit{chooses} to spend and that choice is shaped by institutional quality.

	These findings yield a clear policy prescription, differentiated by institutional maturity. For low-income and lower-middle-income countries at an early stage of institutional development, the priority must be capacity building: strengthening tax administration, improving budget transparency, establishing credible and independent central banking, and broadening access to public services. Until these foundations are set, the austerity channel poses too great a risk, and debt should be deployed with extreme caution and directed toward the institutional investments that will expand future fiscal space. For upper-middle-income and advanced economies that have already achieved a degree of institutional maturity, the evidence supports a more confident use of debt as a counter-cyclical instrument provided that the composition of spending is designed to reach broad segments of the population rather than accruing disproportionately to asset holders and elites.

	For international policy debates, the findings caution decisively against reliance on simple debt-to-GDP thresholds or cross-country averages. The 90 percent threshold popularized by \cite{Reinhart2010}, whatever its merits as a stylized fact, obscures the institutional heterogeneity that our analysis reveals to be the dominant feature of the debt-inequality landscape. Context-sensitive assessments of fiscal space, ones that account for institutional capacity, exchange rate arrangements, and the functional composition of public spending, are essential.

	Future research should extend the analysis along two dimensions. First, incorporating more granular data on the functional composition of public spending would clarify which categories of expenditure drive the equalizing or disequalizing effects of debt across institutional contexts. Second, modeling the interaction between domestic political institutions and external financial constraints would illuminate the mechanisms through which institutional quality mediates the debt-inequality relationship. Doing so would further clarify the conditions under which public debt can support both growth and equity in the global economy, and would provide the empirical foundation for the institutional-first framework we advocate here.



	\bibliographystyle{aea}
	\bibliography{references} % Bibliography file (usually '*.bib' )

\end{document}